\section{System Model and Mission Setting}
\label{sec:system}

\subsection{Architecture}

OrbInspect is organized as a ROS~2 Jazzy workspace with packages for dynamics,
control, guidance, perception, safety, mission management, evaluation,
description, Gazebo visualization, and bringup. The system is designed around
a strict authority boundary. The dynamics node propagates the chaser state
under CW/HCW equations and publishes odometry. The controller tracks reference
states and emits acceleration commands. The safety layer monitors clearance
against station primitives and can filter commanded accelerations. The
coverage layer evaluates which inspection targets are visible. The evaluation
nodes record CSV logs, rosbag data, summaries, and figures. Gazebo Harmonic
subscribes to the resulting state and renders the ISS model and chaser for
visual validation.

\subsection{Reference Frame and State}

The inspection pass is represented in a local-vertical/local-horizontal frame
attached to a nominal circular chief orbit. Over the time horizon considered
here, the inspected structure is treated as fixed in this frame. The chaser
translational state is
\begin{equation}
  x =
  \begin{bmatrix}
    r_x & r_y & r_z & v_x & v_y & v_z
  \end{bmatrix}^{\mathsf{T}},
\end{equation}
where $r \in \R^3$ is relative position and $v \in \R^3$ is relative velocity.
The camera is assumed to point toward the inspected structure or toward a
selected target cluster during the baseline experiments; attitude slew
constraints are left to later work.

\subsection{Relative Dynamics}

The baseline dynamics are
\begin{align}
  \dot r_x &= v_x, &
  \dot r_y &= v_y, &
  \dot r_z &= v_z, \\
  \dot v_x &= 3n^2 r_x + 2n v_y + a_x, &
  \dot v_y &= -2n v_x + a_y, &
  \dot v_z &= -n^2 r_z + a_z,
\end{align}
where $n$ is the chief mean motion and
$u=[a_x,a_y,a_z]^{\mathsf{T}}$ is the commanded acceleration. The discrete
propagation map is denoted
\begin{equation}
  x_{k+1}=f_{\mathrm{CW}}(x_k,u_k,\dt).
\end{equation}
The same model is used in the offline trajectory generator and in the ROS
dynamics node for the baseline validation reported here.

\subsection{Inspection Targets}

The station is represented for planning by a finite set
$\mathcal{T}=\{1,\ldots,N\}$ of target samples. Each target $i$ has a point
$p_i$, outward normal $n_i$, and accumulated dwell time $\tau_i$. A target is
considered inspected when it has remained visible for at least
$\tau_{\min}$. The primary offline benchmark in this paper uses the NASA ISS
mesh backend, area-weighted sampling with 180 target samples, and 357 candidate
viewpoints. The ROS smoke run uses the online coverage configuration, which
reports 142 target samples. Gazebo renders the ISS visual model for
interpretation and trajectory replay; the quantitative coverage and dynamics
claims are computed from the logged offline and ROS data.

\subsection{Mission Termination}

The mission starts from a safe initial relative state and terminates when a
coverage threshold is reached, when the planner exhausts available useful
viewpoints, or when a maximum mission time is reached:
\begin{equation}
  C_K = \frac{1}{N}\sum_{i=1}^{N} y_{i,K},
  \qquad
  y_{i,K}=\mathbf{1}\{\tau_{i,K}\ge \tau_{\min}\}.
\end{equation}
The high-coverage offline benchmark uses a 70\% success threshold and a 98\%
planning stop ratio. The best feasible method reaches 98.33\% final coverage,
as reported in Section~\ref{sec:results}.
